%% cover_letter.tex — C16: SAMBPS Unified 5-Layer Architecture
%% Proceedings of the IEEE — Tutorial/Survey
\documentclass[11pt,a4paper]{letter}
\usepackage[T1]{fontenc}
\usepackage[utf8]{inputenc}
\usepackage[margin=2.5cm]{geometry}
\usepackage{microtype,amsmath}

\signature{Anoop~V.~Eluvathingal\\
           Ph.D.\ Candidate, Dept.\ of Electrical Engineering\\
           Indian Institute of Technology Madras\\
           Chennai 600036, India\\[2pt]
           \texttt{arun.eluvathingal@iitm.ac.in}}

\address{Department of Electrical Engineering\\
         Indian Institute of Technology Madras\\
         Chennai 600036, India}

\begin{document}
\begin{letter}{Editor-in-Chief\\
               \textit{Proceedings of the IEEE}}

\opening{Dear Editor,}

We submit for your consideration the manuscript:

\begin{center}
  \textbf{Self-Adaptive Model-Based Protection for IBR-Penetrated
  Networks: A Unified Five-Layer Framework}
\end{center}

for publication as a \textit{Tutorial and Survey} article in
\textit{Proceedings of the IEEE}.

\textbf{The problem.}
Inverter-based resources (IBRs) --- photovoltaic plants, DFIG wind
turbines, battery storage, and virtual synchronous machines --- are
penetrating distribution and sub-transmission networks at levels that
break the fundamental assumptions of classical protection (overcurrent,
distance, and differential).  IBR current limiters clip fault currents
at 1.1--1.2\,pu, IBR control laws transition modes during sustained
faults, and the aggregate penetration ratio $k_{\rm ibr}$ varies
continuously.  As a result, classical differential relays with fixed
slope settings are losing 7--11 percentage points of zone accuracy
when $k_{\rm ibr} > 0.25$.  There is no unified framework that
combines online IBR estimation, physics-constrained adaptation, and
stability-verified tripping in a single deployable relay algorithm.

\textbf{Contributions.}
This paper presents the \textbf{Self-Adaptive Model-Based Protection
System (SAMBPS)}, a five-layer hierarchical architecture developed
over sixteen theoretical contributions and validated in seven
independent campaigns:

\begin{enumerate}
  \item \textbf{Unified architecture (C16)}: A five-layer framework
        (signal conditioning, LM estimation, adaptive laws, stability
        verification, protection decision) with well-defined layer
        contracts and a single integrating signal --- the
        condition number $\kappa_n$ from the Levenberg--Marquardt
        estimator.

  \item \textbf{$\kappa_n$ as system integrator}: We prove that the
        same threshold $\kappa_n = 30$ is simultaneously the gate
        condition for the adaptive law (C1), the SVD reliability
        criterion for the LM estimator (C9), the CUEP gate (C13),
        and the Stage-2 trip veto (L5).  This triple-role consistency
        is proved in an appendix.

  \item \textbf{Comprehensive theory-to-validation mapping}: All
        sixteen contributions and seven campaigns are allocated to
        specific layers via a contribution map (Fig.~2), showing that
        every layer has at least two independent validation campaigns.

  \item \textbf{Field demonstration}: 95.5\% zone accuracy on
        157 real 33\,kV DSO fault recordings, compared to 89.2\%
        for classical protection on the same dataset.
\end{enumerate}

\textbf{Fit to \textit{Proceedings of the IEEE}.}
\textit{Proceedings of the IEEE} is the primary venue for broad
tutorial surveys that unify a research programme and orient the field
towards future directions.  This paper synthesises five primary
research papers (C1, C9, C13, C14, and seven G-series validations)
into a coherent framework suitable for readers who want to understand
the full SAMBPS system without reading each component paper.  The
architecture is novel, the field validation is the most comprehensive
to date, and the open problems section identifies concrete next steps
for the community.

\textbf{Suggested reviewers.}
\begin{itemize}
  \item Prof.\ Arun G.\ Phadke (Virginia Tech)
        --- wide-area protection, phasor measurement
  \item Prof.\ Ali Hooshyar (University of Toronto)
        --- inverter-based resource protection
  \item Prof.\ Vijay Vittal (Arizona State University)
        --- transient stability, power system dynamics
  \item Prof.\ Petros Ioannou (University of Southern California)
        --- adaptive control, Lyapunov methods
\end{itemize}

This manuscript is not under review elsewhere.
All authors have approved the submission.

\closing{Sincerely,}
\end{letter}
\end{document}
